\documentclass[11pt,a4paper]{article}

% ---------------------------------------------------------------------------
% Paquets
% ---------------------------------------------------------------------------
\usepackage[utf8]{inputenc}
\usepackage[T1]{fontenc}
\usepackage[french]{babel}
\usepackage{lmodern}
\usepackage[margin=2.3cm]{geometry}
\usepackage{amsmath,amssymb}
\usepackage{xcolor}
\usepackage{listings}
\usepackage{booktabs}
\usepackage{tabularx}
\usepackage{array}
\usepackage{float}
\usepackage{caption}
\usepackage{fancyhdr}
\usepackage{algorithm}
\usepackage{algpseudocode}
\usepackage{tikz}
\usetikzlibrary{positioning,arrows.meta,automata,babel}
\usepackage{hyperref}

\hypersetup{colorlinks=true, linkcolor=blue!50!black, urlcolor=blue!50!black}

% ---------------------------------------------------------------------------
% Style du code C++
% ---------------------------------------------------------------------------
\definecolor{fondcode}{RGB}{248,248,250}
\definecolor{motcle}{RGB}{0,0,160}
\definecolor{commentaire}{RGB}{0,120,0}
\definecolor{chaine}{RGB}{170,40,40}

\lstdefinestyle{cpp}{
  language=C++,
  backgroundcolor=\color{fondcode},
  basicstyle=\ttfamily\footnotesize,
  keywordstyle=\color{motcle}\bfseries,
  commentstyle=\color{commentaire}\itshape,
  stringstyle=\color{chaine},
  numbers=left, numberstyle=\tiny\color{gray}, numbersep=6pt,
  frame=single, rulecolor=\color{gray!40},
  breaklines=true, showstringspaces=false, tabsize=4,
  morekeywords={override,nullptr,template,typename},
  literate={é}{{\'e}}1 {è}{{\`e}}1 {ê}{{\^e}}1 {à}{{\`a}}1 {â}{{\^a}}1
           {ô}{{\^o}}1 {î}{{\^i}}1 {ù}{{\`u}}1 {û}{{\^u}}1 {ç}{{\c c}}1
           {É}{{\'E}}1 {ï}{{\"i}}1
}
\lstdefinestyle{sortie}{
  basicstyle=\ttfamily\footnotesize,
  backgroundcolor=\color{black!3},
  frame=leftline, rulecolor=\color{blue!40!black}, numbers=none,
  keywordstyle={}, commentstyle={}, stringstyle={},
  breaklines=true, columns=fullflexible, keepspaces=true
}
\lstset{style=cpp}

% ---------------------------------------------------------------------------
% Pseudo-code en français
% ---------------------------------------------------------------------------
\floatname{algorithm}{Algorithme}
\algrenewcommand\algorithmicfunction{\textbf{fonction}}
\algrenewcommand\algorithmicfor{\textbf{pour}}
\algrenewcommand\algorithmicforall{\textbf{pour tout}}
\algrenewcommand\algorithmicdo{\textbf{faire}}
\algrenewcommand\algorithmicwhile{\textbf{tant que}}
\algrenewcommand\algorithmicif{\textbf{si}}
\algrenewcommand\algorithmicthen{\textbf{alors}}
\algrenewcommand\algorithmicelse{\textbf{sinon}}
\algrenewcommand\algorithmicend{\textbf{fin}}
\algrenewcommand\algorithmicreturn{\textbf{retourner}}
\algrenewcommand\algorithmicrequire{\textbf{Entrée :}}
\algrenewcommand\algorithmicensure{\textbf{Sortie :}}

% Raccourcis
\newcommand{\code}[1]{\texttt{#1}}
\newcommand{\ok}{\textcolor{green!50!black}{\textbf{valide}}}
\newcommand{\ko}{\textcolor{red!70!black}{\textbf{invalide}}}

% En-têtes
\setlength{\headheight}{14pt}
\pagestyle{fancy}
\fancyhf{}
\lhead{TP 1 -- Piles génériques}
\rhead{POO avancée \& Design Patterns}
\cfoot{\thepage}

% ===========================================================================
\begin{document}

% ---------------------------------------------------------------------------
% Page de titre
% ---------------------------------------------------------------------------
\begin{titlepage}
  \centering
  \vspace*{2cm}
  {\large Programmation Orientée Objet Avancée et Design Patterns\par}
  \vspace{1.5cm}
  \rule{\linewidth}{0.6pt}\par\vspace{0.4cm}
  {\Huge\bfseries Rapport de TP 1\par}
  \vspace{0.3cm}
  {\LARGE Pile générique et analyse d'expressions mathématiques\par}
  \vspace{0.4cm}
  \rule{\linewidth}{0.6pt}\par
  \vspace{1.5cm}
  {\large Langage : C++ \quad--\quad Système : Linux\par}
  \vspace{0.4cm}
  {\normalsize Code source : \url{https://github.com/zakariae-bakkari/cpp_ilisi2/tree/main/tp1}\par}
  \vfill
  {\large Réalisé par : \textbf{Zakariae Bakkari}\par}
  \vspace{1cm}
  {\large \today\par}
\end{titlepage}

\tableofcontents
\newpage

% ===========================================================================
\section{Introduction}

La \emph{pile} (\emph{stack}) est une structure de données de type LIFO
(\emph{Last In, First Out}) : le dernier élément ajouté est le premier retiré.
Elle est utilisée dans de nombreux algorithmes : gestion des appels de fonctions,
annulation (\emph{undo}), parcours en profondeur, et surtout
\textbf{analyse syntaxique d'expressions}.

Ce TP comporte trois exercices liés entre eux :
\begin{enumerate}
  \item concevoir une \textbf{pile générique} (\code{template}) avec deux
        implémentations : tableau dynamique et liste chaînée ;
  \item utiliser cette pile pour \textbf{vérifier une expression mathématique}
        (validité, équilibre et ordre des délimiteurs), puis l'évaluer en
        respectant la priorité des opérateurs ;
  \item \textbf{convertir} une expression de la notation infixe
        vers la notation préfixe et inversement.
\end{enumerate}

Pour chaque exercice, nous présentons la \textbf{problématique}, la
\textbf{méthodologie} adoptée et les \textbf{résultats} obtenus.

\paragraph{Environnement de travail.} Le TP est réalisé sous \textbf{Linux}
(Ubuntu, compilateur \code{g++} 15.2, outil \code{make}). Le code des piles
et de l'analyseur d'expressions est compilé sous forme de
\textbf{bibliothèque dynamique} \code{libpile.so}, l'équivalent Linux d'une
DLL Windows, utilisée par les trois programmes de test
(section~\ref{sec:bibliotheque}).

\paragraph{Dépôt GitHub.} L'ensemble du code source, la bibliothèque, les
fichiers d'étapes de compilation et ce rapport sont disponibles à
l'adresse : \url{https://github.com/zakariae-bakkari/cpp_ilisi2/tree/main/tp1}.

\paragraph{Organisation du code.}
\begin{center}
\small
\begin{tabularx}{\linewidth}{l X}
  \toprule
  \textbf{Fichier} & \textbf{Rôle} \\
  \midrule
  \code{include/Pile.h}        & interface abstraite générique \code{Pile<T>} \\
  \code{include/PileTableau.h}, \code{.tpp} & tableau dynamique : déclaration, implémentation \\
  \code{include/PileListe.h}, \code{.tpp}   & liste chaînée : déclaration, implémentation \\
  \code{include/Expression.h}, \code{.tpp}  & analyse lexicale, vérification, évaluation, conversions \\
  \code{src/PileTableau.cpp}   & instanciations de \code{PileTableau<T>} (bibliothèque) \\
  \code{src/PileListe.cpp}     & instanciations de \code{PileListe<T>} (bibliothèque) \\
  \code{src/Expression.cpp}    & fonctions utilitaires et \code{Analyseur} (bibliothèque) \\
  \code{src/test\_pile.cpp}       & tests de l'exercice 1 \\
  \code{src/test\_expression.cpp} & tests de l'exercice 2 \\
  \code{src/test\_conversion.cpp} & tests de l'exercice 3 \\
  \code{Makefile}              & compilation de \code{libpile.so} et des tests \\
  \code{compilation\_steps/}   & fichiers générés à chaque étape de compilation \\
  \code{guide/}                & guide PDF : bibliothèques \code{.so} / \code{.a}, fichiers \code{.tpp} \\
  \bottomrule
\end{tabularx}
\end{center}

\noindent Pour la lisibilité, les extraits de code des sections 2 à 4
présentent les méthodes écrites à l'intérieur des classes ; dans le projet,
leur code se trouve dans les fichiers \code{.tpp}
(voir section~\ref{sec:bibliotheque}).

% ===========================================================================
\section{Exercice 1 : Pile générique}

\subsection{Problématique}

Une pile écrite pour un type fixe (par exemple \code{int}) doit être
réécrite pour chaque nouveau type (\code{char}, \code{string},
\code{double}\dots), ce qui duplique le code et multiplie les erreurs.
De plus, il existe plusieurs représentations possibles d'une pile, chacune
avec ses avantages : il faut pouvoir en changer \emph{sans modifier le code
qui l'utilise}.

Les questions posées sont donc :
\begin{itemize}
  \item Comment écrire une pile \textbf{indépendante du type} des éléments ?
  \item Comment offrir \textbf{deux représentations} (tableau dynamique et
        liste chaînée) derrière un \textbf{même contrat} ?
  \item Comment garantir une \textbf{gestion mémoire correcte}
        (pas de fuite, copie profonde, pile vide) ?
  \item Quelle représentation est la plus \textbf{performante} ?
\end{itemize}

\subsection{Méthodologie}

\subsubsection{Conception}

Nous séparons le \emph{contrat} de l'\emph{implémentation} :
\begin{itemize}
  \item une classe abstraite générique \code{Pile<T>} déclare les opérations
        (\code{empiler}, \code{depiler}, \code{getSommet}, \code{estVide},
        \code{taille}) sous forme de méthodes virtuelles pures ;
  \item deux classes génériques concrètes, \code{PileTableau<T>} et
        \code{PileListe<T>}, héritent de cette interface ;
  \item le code client (fonctions de test, analyseur d'expressions) ne dépend
        que de l'interface ou du paramètre de modèle, ce qui permet
        d'interchanger les implémentations.
\end{itemize}

Cette conception combine les deux formes de polymorphisme du C++ :
le polymorphisme \textbf{paramétrique} (templates, choix du type \code{T})
et le polymorphisme \textbf{d'inclusion} (héritage et méthodes virtuelles,
choix de la représentation).

\begin{figure}[H]
\centering
\begin{tikzpicture}[
    classe/.style={draw, fill=blue!4, inner sep=0pt, font=\small},
    herite/.style={-{Triangle[open,length=3mm,width=3mm]}, thick},
    utilise/.style={dashed, -{Stealth[length=2.5mm]}}
  ]
  \node[classe] (pile) {
    \begin{tabular}{l}
      \multicolumn{1}{c}{\guillemotleft interface\guillemotright} \\
      \multicolumn{1}{c}{\textbf{\textit{Pile<T>}}} \\ \hline
      \code{+ empiler(x : const T\&) : void} \\
      \code{+ depiler() : T} \\
      \code{+ getSommet() : const T\&} \\
      \code{+ estVide() : bool} \\
      \code{+ taille() : int}
    \end{tabular}};

  \node[classe, below left=1.4cm and -1.2cm of pile] (tab) {
    \begin{tabular}{l}
      \multicolumn{1}{c}{\textbf{PileTableau<T>}} \\ \hline
      \code{- tab : T*} \\
      \code{- capacite : int} \\
      \code{- nb : int} \\ \hline
      \code{- agrandir() : void} \\
      \code{+ getCapacite() : int} \\
      \code{+ afficher() : void}
    \end{tabular}};

  \node[classe, below right=1.4cm and -1.2cm of pile] (lst) {
    \begin{tabular}{l}
      \multicolumn{1}{c}{\textbf{PileListe<T>}} \\ \hline
      \code{- tete : Noeud*} \\
      \code{- nb : int} \\ \hline
      \code{- vider() : void} \\
      \code{- copier(p) : void} \\
      \code{+ afficher() : void}
    \end{tabular}};

  \node[classe, right=1.2cm of pile] (ana) {
    \begin{tabular}{l}
      \multicolumn{1}{c}{\textbf{Analyseur<P>}} \\ \hline
      \code{+ verifier(expr)} \\
      \code{+ evaluer(expr)} \\
      \code{+ versPrefixe(expr)} \\
      \code{+ prefixeVersInfixe(pre)}
    \end{tabular}};

  \draw[thick] (tab.north) |- ([yshift=-0.7cm]pile.south);
  \draw[thick] (lst.north) |- ([yshift=-0.7cm]pile.south);
  \draw[herite] ([yshift=-0.7cm]pile.south) -- (pile.south);
  \draw[utilise] (ana.west) -- node[above, font=\scriptsize]{utilise} (pile.east);
\end{tikzpicture}
\caption{Diagramme de classes simplifié}
\end{figure}

\subsubsection{Interface générique}

\begin{lstlisting}[caption={Interface \code{Pile<T>} (\code{Pile.h})}]
template <typename T>
class Pile {
public:
    virtual ~Pile() {}

    virtual void empiler(const T& x) = 0;       // ajoute x au sommet
    virtual T depiler() = 0;                    // retire et renvoie le sommet
    virtual const T& getSommet() const = 0;     // consulte le sommet
    virtual bool estVide() const = 0;
    virtual int taille() const = 0;
};
\end{lstlisting}

Le destructeur virtuel garantit que la destruction d'une pile concrète via
un pointeur \code{Pile<T>*} libère correctement la mémoire.

\subsubsection{Implémentation par tableau dynamique}

Les éléments sont rangés dans un tableau alloué sur le tas ; le sommet est
la case \code{tab[nb-1]}. Lorsque le tableau est plein, sa capacité est
\textbf{doublée} : une copie en $O(n)$ n'a lieu que $\log_2 n$ fois, ce qui
donne un coût \textbf{amorti} en $O(1)$ par empilement.

\begin{figure}[H]
\centering
\begin{tikzpicture}[
    case/.style={draw, minimum width=0.9cm, minimum height=0.7cm, font=\small},
    libre/.style={case, fill=gray!15}
  ]
  \foreach \i/\v in {0/10,1/20,2/30,3/40,4/50,5/60}
    \node[case, fill=blue!8] (t\i) at (\i*0.9,0) {\v};
  \foreach \i in {6,7}
    \node[libre] (t\i) at (\i*0.9,0) {};
  \foreach \i in {0,...,7}
    \node[font=\scriptsize, gray, below=1pt of t\i] {\i};
  \draw[{Stealth[length=2.5mm]}-] (t5.north) -- ++(0,0.7)
    node[above, font=\small] {sommet = \code{tab[nb-1]}};
  \node[font=\small, right=0.4cm of t7, align=left]
    {\code{nb = 6}\\ \code{capacite = 8}};
\end{tikzpicture}
\caption{Représentation de la pile par tableau dynamique}
\end{figure}

\begin{lstlisting}[caption={Méthodes essentielles de \code{PileTableau<T>}}]
// double la capacité quand le tableau est plein : coût amorti O(1)
void agrandir() {
    int nouvelleCapacite = capacite * 2;
    T* nouveau = new T[nouvelleCapacite];
    for (int i = 0; i < nb; i++)
        nouveau[i] = tab[i];
    delete[] tab;
    tab = nouveau;
    capacite = nouvelleCapacite;
}

void empiler(const T& x) override {
    if (nb == capacite)
        agrandir();
    tab[nb++] = x;
}

T depiler() override {
    if (estVide())
        throw std::underflow_error("Pile vide");
    return tab[--nb];
}
\end{lstlisting}

\subsubsection{Implémentation par liste chaînée}

Chaque élément est stocké dans un \code{Noeud} alloué individuellement ;
le sommet est la \textbf{tête} de la liste, ce qui rend l'empilement et le
dépilement en $O(1)$ sans jamais recopier les données.

\begin{figure}[H]
\centering
\begin{tikzpicture}[
    val/.style={draw, fill=blue!8, minimum width=0.9cm, minimum height=0.7cm, font=\small},
    ptr/.style={draw, minimum width=0.5cm, minimum height=0.7cm}
  ]
  \node (tete) {\code{tete}};
  \node[val, right=0.8cm of tete] (v1) {60};
  \node[ptr, right=0pt of v1] (p1) {};
  \node[val, right=0.8cm of p1] (v2) {50};
  \node[ptr, right=0pt of v2] (p2) {};
  \node[val, right=0.8cm of p2] (v3) {40};
  \node[ptr, right=0pt of v3] (p3) {};
  \node[right=0.3cm of p3] (dots) {$\cdots$};
  \node[val, right=0.3cm of dots] (v4) {10};
  \node[ptr, right=0pt of v4] (p4) {};
  \node[right=0.6cm of p4, font=\small] (nul) {\code{nullptr}};

  \draw[-{Stealth[length=2.5mm]}] (tete) -- (v1);
  \foreach \src/\dst in {p1/v2, p2/v3, p3/dots, p4/nul} {
    \fill (\src.center) circle (1.5pt);
    \draw[-{Stealth[length=2.5mm]}] (\src.center) -- (\dst.west);
  }
\end{tikzpicture}
\caption{Représentation de la pile par liste chaînée (sommet = tête)}
\end{figure}

\begin{lstlisting}[caption={Structure et méthodes essentielles de \code{PileListe<T>}}]
struct Noeud {
    T valeur;
    Noeud* suivant;
    Noeud(const T& v, Noeud* s = nullptr) : valeur(v), suivant(s) {}
};

void empiler(const T& x) override {
    tete = new Noeud(x, tete);   // insertion en tête
    nb++;
}

T depiler() override {
    if (estVide())
        throw std::underflow_error("Pile vide");
    Noeud* tmp = tete;
    T v = tmp->valeur;
    tete = tete->suivant;
    delete tmp;
    nb--;
    return v;
}
\end{lstlisting}

\subsubsection{Gestion de la mémoire}

Les deux classes possèdent de la mémoire dynamique ; nous appliquons donc la
\textbf{règle de trois} : destructeur, constructeur de copie et opérateur
d'affectation effectuent une \textbf{copie profonde}. Pour
\code{PileTableau}, l'opérateur d'affectation alloue le nouveau tableau
\emph{avant} de libérer l'ancien afin de rester cohérent si \code{new}
échoue. Toute opération sur une pile vide lève l'exception
\code{std::underflow\_error}.

\subsubsection{Complexité}

\begin{table}[H]
\centering
\begin{tabular}{lcc}
  \toprule
  \textbf{Opération} & \textbf{PileTableau} & \textbf{PileListe} \\
  \midrule
  \code{empiler}    & $O(1)$ amorti ($O(n)$ lors d'un agrandissement) & $O(1)$ \\
  \code{depiler}    & $O(1)$ & $O(1)$ \\
  \code{getSommet}  & $O(1)$ & $O(1)$ \\
  copie             & $O(n)$ & $O(n)$ \\
  mémoire           & entre $n$ et $2n$ cases & $n$ n\oe{}uds (valeur + pointeur) \\
  \bottomrule
\end{tabular}
\caption{Comparaison théorique des deux implémentations}
\end{table}

\subsection{Résultats}

Le programme \code{test\_pile} valide chaque aspect de la conception :

\begin{lstlisting}[style=sortie, caption={Extrait de la sortie de \code{test\_pile}}]
=== 1. PileTableau<int> ===
  empiler(10) -> taille=1 capacite=2
  empiler(20) -> taille=2 capacite=2
  empiler(30) -> taille=3 capacite=4
  empiler(40) -> taille=4 capacite=4
  empiler(50) -> taille=5 capacite=8
  empiler(60) -> taille=6 capacite=8
  contenu : [ 60 50 40 30 20 10 ]
  sommet  : 60
  depiler : 60

=== 2. PileListe<string> ===
  contenu : [ gamma beta alpha ]

=== 3. PileTableau<Point> (type utilisateur) ===
  contenu : [ (3,4) (1,2) ]

=== 4. Polymorphisme via Pile<T>& ===
PileTableau<char>
  depilement : E L I P
PileListe<char>
  depilement : E L I P

=== 5. Copie profonde ===
  original : [ 2 1 ]
  copie    : [ 99 2 1 ]

=== 6. Exceptions (pile vide) ===
  PileTableau : Pile vide
  PileListe   : Pile vide
\end{lstlisting}

\paragraph{Interprétation.}
\begin{itemize}
  \item \textbf{Généricité :} la même classe fonctionne avec \code{int},
        \code{string}, \code{char} et un type utilisateur \code{Point}.
  \item \textbf{Agrandissement :} la capacité double (2, 4, 8) exactement
        lorsque le tableau est plein.
  \item \textbf{Polymorphisme :} la fonction \code{viderEtAfficher(Pile<T>\&)}
        traite les deux implémentations de façon identique ; empiler
        \og PILE \fg{} puis tout dépiler donne bien \og ELIP \fg{} (ordre LIFO).
  \item \textbf{Copie profonde :} modifier la copie ne modifie pas l'original.
  \item \textbf{Robustesse :} le dépilement d'une pile vide est intercepté par
        une exception.
\end{itemize}

\paragraph{Performances.} Temps mesurés pour empiler puis dépiler $n$
entiers (compilation \code{-O2}) :

\begin{table}[H]
\centering
\begin{tabular}{rrrr}
  \toprule
  $n$ & \textbf{PileTableau} (ms) & \textbf{PileListe} (ms) & \textbf{Rapport} \\
  \midrule
  $10^5$ & 0{,}30  & 3{,}03   & $\approx 10\times$ \\
  $10^6$ & 2{,}14  & 36{,}18  & $\approx 17\times$ \\
  $10^7$ & 32{,}24 & 320{,}12 & $\approx 10\times$ \\
  \bottomrule
\end{tabular}
\caption{Comparaison expérimentale des deux implémentations}
\end{table}

Les deux implémentations ont un temps linéaire en $n$, mais le tableau
dynamique est environ \textbf{dix fois plus rapide}. La liste effectue une
allocation (\code{new}) et une libération (\code{delete}) par élément,
tandis que le tableau ne réalloue que $\log_2 n$ fois ($\approx 22$ fois
pour $n = 10^7$) et stocke ses éléments de façon contiguë, ce qui exploite
le cache du processeur. La liste reste intéressante lorsque l'on veut éviter
les recopies (éléments coûteux à copier) ou garantir un temps constant
\emph{dans le pire cas} pour chaque opération.

% ===========================================================================
\section{Exercice 2 : Vérification d'une expression mathématique}

\subsection{Problématique}

Avant d'évaluer une expression comme
$\{\,[\,(1+2)\times 3\,]-4\,\}\,/\,5$, il faut s'assurer qu'elle est
\textbf{bien formée}. Plusieurs types d'erreurs doivent être détectés :
\begin{enumerate}
  \item \textbf{caractères interdits} : \code{3 \# 4} ;
  \item \textbf{délimiteurs non équilibrés} : \code{(3 + 4}, \code{3 + 4)} ;
  \item \textbf{délimiteurs mal appariés} : \code{(3 + 4]} ;
  \item \textbf{ordre des délimiteurs non respecté} : par convention
        mathématique, les accolades contiennent les crochets qui contiennent
        les parenthèses, $\{\,[\,(\;)\,]\,\}$ ; ainsi \code{( [ 1 + 2 ] )}
        est refusé ;
  \item \textbf{erreurs de syntaxe} entre opérandes et opérateurs :
        \code{3 + * 4}, \code{3 4 + 5}, \code{(3 + )}, \code{()}, \code{3 +}.
\end{enumerate}
Une fois l'expression validée, il faut pouvoir l'\textbf{évaluer} en
respectant la \textbf{priorité} (\code{\textasciicircum} $>$ \code{* /} $>$ \code{+ -}) et
l'\textbf{associativité} des opérateurs. La difficulté est que
l'imbrication des délimiteurs est arbitrairement profonde : un simple
compteur ne suffit pas à vérifier l'appariement, il faut \emph{mémoriser
l'ordre} des ouvertures, d'où l'usage d'une pile.

\subsection{Méthodologie}

\subsubsection{Analyse lexicale}

La fonction \code{decouper} transforme la chaîne en une suite de
\textbf{jetons} (\code{Token}) : nombre (entier ou décimal), variable
(identificateur), opérateur (\code{+ - * / \textasciicircum}), délimiteur
ouvrant (\code{( [ \{}) ou fermant (\code{) ] \}}). Chaque jeton mémorise sa
\textbf{position} dans la chaîne afin de localiser précisément les erreurs.
Les espaces sont ignorés et tout autre caractère déclenche une erreur.

\subsubsection{Automate de vérification}

La vérification combine deux mécanismes :
\begin{itemize}
  \item un \textbf{automate à deux états} qui contrôle l'alternance
        opérande/opérateur ;
  \item une \textbf{pile générique de jetons} (\code{P<Token>}) qui mémorise
        les délimiteurs ouverts et non encore fermés.
\end{itemize}

\begin{figure}[H]
\centering
\begin{tikzpicture}[
    ->, >={Stealth[length=2.5mm]}, node distance=6cm, auto,
    every state/.style={minimum size=2.4cm, align=center, font=\small},
    lab/.style={font=\footnotesize, align=center}
  ]
  \node[state, initial, initial text={début}] (O) {Attend\\ opérande};
  \node[state, accepting, right=of O] (P) {Attend\\ opérateur};

  \path
    (O) edge[bend left=20]  node[lab]{nombre / variable} (P)
    (P) edge[bend left=20]  node[lab]{opérateur \code{+ - * / \textasciicircum}} (O)
    (O) edge[loop above]    node[lab]{\code{( [ \{}\\ empiler (si ordre respecté)} (O)
    (P) edge[loop above]    node[lab]{\code{) ] \}}\\ dépiler et comparer} (P);
\end{tikzpicture}
\caption{Automate de vérification. Toute autre transition est une erreur ;
l'expression est acceptée si l'on termine dans l'état \og attend opérateur \fg{}
avec une pile vide.}
\end{figure}

Le rang d'un délimiteur est défini par $\text{rang}(\,(\,)=1$,
$\text{rang}(\,[\,)=2$, $\text{rang}(\,\{\,)=3$. Un délimiteur ouvrant
n'est accepté que si son rang est \textbf{inférieur ou égal} à celui du
sommet de la pile : on peut écrire \code{((a))} ou \code{\{[(a)]\}}, mais
pas \code{([a])}.

\begin{algorithm}[H]
\caption{Vérification d'une expression infixe}
\begin{algorithmic}[1]
\Require suite de jetons $t_1 \dots t_k$
\Ensure valide, ou message d'erreur et position
\State $pile \gets$ pile vide ; $attendOperande \gets$ vrai
\ForAll{jeton $t$}
  \If{$t$ est un nombre ou une variable}
    \If{\textbf{non} $attendOperande$} \Return erreur \og opérateur manquant \fg{} \EndIf
    \State $attendOperande \gets$ faux
  \ElsIf{$t$ est un opérateur}
    \If{$attendOperande$} \Return erreur \og opérande manquant \fg{} \EndIf
    \State $attendOperande \gets$ vrai
  \ElsIf{$t$ est un délimiteur ouvrant}
    \If{\textbf{non} $attendOperande$} \Return erreur \og opérateur manquant \fg{} \EndIf
    \If{$pile$ non vide \textbf{et} rang$(t) >$ rang(sommet$(pile)$)} \Return erreur d'ordre \EndIf
    \State empiler$(pile, t)$
  \Else \Comment{délimiteur fermant}
    \If{$pile$ vide} \Return erreur \og fermant sans ouvrant \fg{} \EndIf
    \If{$attendOperande$} \Return erreur \og opérande manquant \fg{} \EndIf
    \State $o \gets$ dépiler$(pile)$
    \If{$o$ n'est pas l'ouvrant associé à $t$} \Return erreur d'appariement \EndIf
  \EndIf
\EndFor
\If{$pile$ non vide} \Return erreur \og délimiteur jamais fermé \fg{} \EndIf
\If{$attendOperande$} \Return erreur \og expression incomplète \fg{} \EndIf
\State \Return valide
\end{algorithmic}
\end{algorithm}

\subsubsection{Généricité vis-à-vis de l'implémentation de la pile}

L'analyseur est une classe paramétrée par la \emph{classe de pile}
elle-même (paramètre de modèle de modèle) : \code{Analyseur<PileTableau>}
et \code{Analyseur<PileListe>} exécutent le même algorithme avec
une représentation différente.

\begin{lstlisting}[caption={C\oe{}ur de la méthode \code{verifier}}]
template <template <typename> class P>
class Analyseur {
public:
    static Resultat verifier(const std::string& expr) {
        // ... découpage en jetons, expression vide ...
        P<Token> delimiteurs;
        bool attendOperande = true;   // état de l'automate

        for (const Token& t : tokens) {
            switch (t.type) {
                case NOMBRE:
                case VARIABLE:
                    if (!attendOperande)
                        return {false, "Operande '" + t.texte + "' inattendu (operateur manquant)", t.position};
                    attendOperande = false;
                    break;
                case OPERATEUR:
                    if (attendOperande)
                        return {false, "Operateur '" + t.texte + "' inattendu (operande manquant)", t.position};
                    attendOperande = true;
                    break;
                case OUVRANTE:
                    if (!attendOperande)
                        return {false, "'" + t.texte + "' inattendu (operateur manquant)", t.position};
                    if (!delimiteurs.estVide()
                        && rangDelimiteur(t.texte[0]) > rangDelimiteur(delimiteurs.getSommet().texte[0]))
                        return {false, "Ordre non respecte : ...", t.position};
                    delimiteurs.empiler(t);
                    break;
                case FERMANTE: {
                    if (delimiteurs.estVide())
                        return {false, "'" + t.texte + "' sans delimiteur ouvrant", t.position};
                    if (attendOperande)
                        return {false, "Operande manquant avant '" + t.texte + "'", t.position};
                    Token ouvrant = delimiteurs.depiler();
                    if (ouvrant.texte[0] != ouvranteAssociee(t.texte[0]))
                        return {false, "'" + t.texte + "' ne ferme pas '" + ouvrant.texte + "'", t.position};
                    break;
                }
            }
        }
        if (!delimiteurs.estVide())
            return {false, "'" + delimiteurs.getSommet().texte + "' jamais ferme", ...};
        if (attendOperande)
            return {false, "Expression incomplete (operande manquant a la fin)", ...};
        return {true, "Expression valide", -1};
    }
};
\end{lstlisting}

\subsubsection{Trace d'exécution}

Pour l'expression \code{\{ [ (1 + 2) * 3 ] - 4 \} / 5} (sommet de la pile à
droite ; O = attend opérande, P = attend opérateur) :

\begin{table}[H]
\centering
\small
\begin{tabular}{clllc}
  \toprule
  \textbf{Étape} & \textbf{Jeton} & \textbf{Action} & \textbf{Pile} & \textbf{État} \\
  \midrule
  1  & \code{\{} & empiler                        & \code{\{}       & O \\
  2  & \code{[}  & rang 2 $\le$ 3 : empiler       & \code{\{ [}     & O \\
  3  & \code{(}  & rang 1 $\le$ 2 : empiler       & \code{\{ [ (}   & O \\
  4  & \code{1}  & opérande                       & \code{\{ [ (}   & P \\
  5  & \code{+}  & opérateur                      & \code{\{ [ (}   & O \\
  6  & \code{2}  & opérande                       & \code{\{ [ (}   & P \\
  7  & \code{)}  & dépiler \code{(} : apparié     & \code{\{ [}     & P \\
  8  & \code{*}  & opérateur                      & \code{\{ [}     & O \\
  9  & \code{3}  & opérande                       & \code{\{ [}     & P \\
  10 & \code{]}  & dépiler \code{[} : apparié     & \code{\{}       & P \\
  11 & \code{-}  & opérateur                      & \code{\{}       & O \\
  12 & \code{4}  & opérande                       & \code{\{}       & P \\
  13 & \code{\}} & dépiler \code{\{} : apparié    & (vide)          & P \\
  14 & \code{/}  & opérateur                      & (vide)          & O \\
  15 & \code{5}  & opérande                       & (vide)          & P \\
  \midrule
  \multicolumn{5}{c}{Fin : pile vide et état P $\Rightarrow$ expression \ok} \\
  \bottomrule
\end{tabular}
\caption{Trace de la vérification}
\end{table}

\subsubsection{Évaluation : algorithme à deux piles}

Pour évaluer une expression valide, nous utilisons l'algorithme à deux piles
de Dijkstra : une pile d'\textbf{opérateurs} (\code{P<char>}) et une pile de
\textbf{valeurs} (\code{P<V>}).
\begin{itemize}
  \item un opérande est empilé sur la pile des valeurs ;
  \item un délimiteur ouvrant est empilé sur la pile des opérateurs ;
  \item un délimiteur fermant déclenche des \emph{réductions} jusqu'à
        l'ouvrant correspondant ;
  \item un opérateur $op$ déclenche des réductions tant que le sommet est
        un opérateur de priorité \emph{supérieure}, ou \emph{égale} si $op$
        est associatif à gauche, puis $op$ est empilé ;
  \item une \textbf{réduction} dépile un opérateur et deux valeurs $a, b$
        et empile \code{combiner(op, a, b)}.
\end{itemize}

L'algorithme est écrit \textbf{une seule fois}, de façon générique : le type
des valeurs \code{V} et la fonction \code{combiner} sont des paramètres.
Avec \code{V = double} et \code{combiner} = calcul, on obtient
l'\textbf{évaluation} ; avec \code{V = string} et \code{combiner} =
concaténation, on obtient les \textbf{conversions} de l'exercice 3.

\begin{lstlisting}[caption={Parcours générique d'une expression infixe}]
template <typename V>
static V parcourirInfixe(const std::string& expr,
                         std::function<V(const Token&)> feuille,
                         std::function<V(char, const V&, const V&)> combiner) {
    Resultat r = verifier(expr);
    if (!r.valide)
        throw ErreurExpression(r.message, r.position);

    P<char> operateurs;
    P<V> valeurs;

    // dépile un opérateur et ses deux opérandes, empile le résultat
    auto reduire = [&]() {
        char op = operateurs.depiler();
        V b = valeurs.depiler();
        V a = valeurs.depiler();
        valeurs.empiler(combiner(op, a, b));
    };

    for (const Token& t : decouper(expr)) {
        switch (t.type) {
            case NOMBRE:
            case VARIABLE:
                valeurs.empiler(feuille(t));
                break;
            case OUVRANTE:
                operateurs.empiler(t.texte[0]);
                break;
            case FERMANTE:
                while (!estOuvrante(operateurs.getSommet()))
                    reduire();
                operateurs.depiler();   // retire le délimiteur ouvrant
                break;
            case OPERATEUR: {
                char op = t.texte[0];
                while (!operateurs.estVide() && estOperateur(operateurs.getSommet())) {
                    char haut = operateurs.getSommet();
                    if (priorite(haut) > priorite(op)
                        || (priorite(haut) == priorite(op) && !associatifDroite(op)))
                        reduire();
                    else
                        break;
                }
                operateurs.empiler(op);
                break;
            }
        }
    }
    while (!operateurs.estVide())
        reduire();
    return valeurs.depiler();
}

static double evaluer(const std::string& expr) {
    return parcourirInfixe<double>(expr,
        [](const Token& t) -> double {
            if (t.type == VARIABLE)
                throw std::invalid_argument("Variable '" + t.texte + "' non evaluable");
            return std::stod(t.texte);
        },
        [](char op, const double& a, const double& b) { return appliquer(op, a, b); });
}
\end{lstlisting}

\subsection{Résultats}

\subsubsection{Vérification}

Le programme \code{test\_expression} soumet 20 expressions à l'analyseur.
Chaque erreur est localisée par un curseur \code{\textasciicircum} :

\begin{lstlisting}[style=sortie, caption={Extrait de la sortie de \code{test\_expression}}]
  "{ [ ( 1 + 2 ] ) }"
               ^
   -> INVALIDE (pos 12) : ']' ne ferme pas '(' (position 4)
  "( [ 1 + 2 ] )"
     ^
   -> INVALIDE (pos 2) : Ordre non respecte : '[' a l'interieur de '('
  "3 + * 4"
       ^
   -> INVALIDE (pos 4) : Operateur '*' inattendu (operande manquant)
\end{lstlisting}

\begin{table}[H]
\centering
\small
\begin{tabularx}{\linewidth}{>{\ttfamily}l c X}
  \toprule
  \normalfont\textbf{Expression} & \textbf{Verdict} & \textbf{Diagnostic (position)} \\
  \midrule
  3 + 4 * 2                     & \ok & -- \\
  (a + b) * c                   & \ok & -- \\
  \{ [ (1 + 2) * 3 ] - 4 \} / 5 & \ok & -- \\
  {}[(x - y) * (x + y)] \textasciicircum{} 2 & \ok & -- \\
  ((2))                         & \ok & -- \\
  \midrule
  (3 + 4                        & \ko & \code{(} jamais fermé (0) \\
  3 + 4)                        & \ko & \code{)} sans délimiteur ouvrant (5) \\
  (3 + 4]                       & \ko & \code{]} ne ferme pas \code{(} (6) \\
  \{ [ ( 1 + 2 ] ) \}           & \ko & \code{]} ne ferme pas \code{(} (12) \\
  \midrule
  ( [ 1 + 2 ] )                 & \ko & ordre : \code{[} à l'intérieur de \code{(} (2) \\
  {}[ \{ a \} ]                   & \ko & ordre : \code{\{} à l'intérieur de \code{[} (2) \\
  \midrule
  3 + * 4                       & \ko & opérateur \code{*} inattendu (4) \\
  3 4 + 5                       & \ko & opérande \code{4} inattendu (2) \\
  (3 + )                        & \ko & opérande manquant avant \code{)} (5) \\
  ()                            & \ko & opérande manquant avant \code{)} (1) \\
  2 (3 + 4)                     & \ko & \code{(} inattendu, opérateur manquant (2) \\
  * 3                           & \ko & opérateur \code{*} inattendu (0) \\
  3 +                           & \ko & expression incomplète (2) \\
  3 \# 4                        & \ko & caractère non autorisé \code{\#} (2) \\
  \normalfont(chaîne vide)      & \ko & expression vide (0) \\
  \bottomrule
\end{tabularx}
\caption{Résultats de la vérification (identiques avec \code{PileTableau}
et \code{PileListe} : 0 différence)}
\end{table}

Toutes les catégories d'erreurs de la problématique sont détectées, avec
un message explicite et la position exacte du jeton fautif. Les deux
implémentations de la pile produisent des verdicts identiques,
ce qui confirme qu'elles respectent le même contrat.

\subsubsection{Évaluation}

\begin{table}[H]
\centering
\begin{tabular}{>{\ttfamily}l r l}
  \toprule
  \normalfont\textbf{Expression} & \textbf{Résultat} & \textbf{Propriété vérifiée} \\
  \midrule
  3 + 4 * 2                         & 11  & priorité de $\times$ sur $+$ \\
  (3 + 4) * 2                       & 14  & parenthèses \\
  \{ [ (1 + 2) * 3 ] - 4 \} / 5     & 1   & délimiteurs imbriqués \\
  10 - 4 - 3                        & 3   & associativité à gauche \\
  100 / 10 / 5                      & 2   & associativité à gauche \\
  2 \textasciicircum{} 3 \textasciicircum{} 2 & 512 & associativité à droite : $2^{(3^2)}$ \\
  2.5 * 4                           & 10  & nombres décimaux \\
  1 / (2 - 2)                       & erreur & division par zéro \\
  x + 1                             & erreur & variable non évaluable \\
  \bottomrule
\end{tabular}
\caption{Résultats de l'évaluation}
\end{table}

% ===========================================================================
\section{Exercice 3 : Conversion infixe \texorpdfstring{$\leftrightarrow$}{<->} préfixe}

\subsection{Problématique}

En notation \textbf{infixe} ($a + b$), l'opérateur est placé entre ses
opérandes : elle est naturelle pour l'humain mais nécessite des parenthèses
et des règles de priorité. En notation \textbf{préfixe} ou polonaise
($+\ a\ b$), l'opérateur précède ses opérandes : l'expression est
\textbf{non ambiguë sans parenthèses} et s'évalue simplement avec une pile.
Le problème est de passer d'une notation à l'autre :
\begin{itemize}
  \item \textbf{infixe $\to$ préfixe} en respectant la priorité et
        l'associativité (\code{a - b - c} $\to$ \code{- - a b c} mais
        \code{a \textasciicircum{} b \textasciicircum{} c} $\to$
        \code{\textasciicircum{} a \textasciicircum{} b c}) ;
  \item \textbf{préfixe $\to$ infixe} en réintroduisant les parenthèses
        nécessaires, et en détectant les expressions préfixes mal formées.
\end{itemize}

\subsection{Méthodologie}

\subsubsection{Infixe vers préfixe}

Nous réutilisons directement le parcours générique à deux piles de
l'exercice 2, avec \code{V = string} : chaque réduction construit la chaîne
\og \code{op a b} \fg{}. La conversion vers la notation postfixe s'obtient de
la même façon avec \og \code{a b op} \fg{}.

\begin{lstlisting}[caption={Conversions obtenues par spécialisation du parcours générique}]
// Infixe -> préfixe (notation polonaise) : "op a b"
static std::string versPrefixe(const std::string& expr) {
    return parcourirInfixe<std::string>(expr,
        [](const Token& t) { return t.texte; },
        [](char op, const std::string& a, const std::string& b) {
            return std::string(1, op) + " " + a + " " + b;
        });
}

// Infixe -> postfixe (notation polonaise inverse) : "a b op"
static std::string versPostfixe(const std::string& expr) {
    return parcourirInfixe<std::string>(expr,
        [](const Token& t) { return t.texte; },
        [](char op, const std::string& a, const std::string& b) {
            return a + " " + b + " " + std::string(1, op);
        });
}
\end{lstlisting}

\begin{table}[H]
\centering
\small
\begin{tabular}{clll}
  \toprule
  \textbf{Jeton} & \textbf{Action} & \textbf{Pile opérateurs} & \textbf{Pile valeurs} \\
  \midrule
  \code{a} & empiler valeur                          & --              & \code{a} \\
  \code{+} & pile vide : empiler                     & \code{+}        & \code{a} \\
  \code{b} & empiler valeur                          & \code{+}        & \code{a | b} \\
  \code{*} & priorité(\code{+}) $<$ priorité(\code{*}) : empiler & \code{+ *} & \code{a | b} \\
  \code{c} & empiler valeur                          & \code{+ *}      & \code{a | b | c} \\
  fin      & réduire \code{*}                        & \code{+}        & \code{a | * b c} \\
  fin      & réduire \code{+}                        & --              & \code{+ a * b c} \\
  \bottomrule
\end{tabular}
\caption{Trace de la conversion de \code{a + b * c} en préfixe}
\end{table}

\subsubsection{Préfixe vers infixe}

Une expression préfixe se lit \textbf{de droite à gauche} avec une seule
pile :
\begin{itemize}
  \item un opérande est empilé ;
  \item un opérateur dépile son premier opérande $a$ puis son second $b$,
        et empile \og \code{(a op b)} \fg{} ;
  \item à la fin, la pile doit contenir exactement un élément ; sinon
        l'expression est mal formée (opérandes ou opérateurs manquants).
\end{itemize}
Ce parcours est lui aussi générique : avec \code{V = double}, il
\textbf{évalue} directement l'expression préfixe.

\begin{lstlisting}[caption={Parcours générique d'une expression préfixe}]
template <typename V>
static V parcourirPrefixe(const std::string& prefixe,
                          std::function<V(const std::string&)> feuille,
                          std::function<V(char, const V&, const V&)> combiner) {
    // ... découpage de la chaîne en jetons séparés par des espaces ...
    P<V> valeurs;
    for (int i = (int)jetons.size() - 1; i >= 0; i--) {
        const std::string& s = jetons[i];
        if (s.size() == 1 && estOperateur(s[0])) {
            if (valeurs.taille() < 2)
                throw std::invalid_argument("Operandes manquants pour '" + s + "'");
            V a = valeurs.depiler();   // premier opérande
            V b = valeurs.depiler();   // second opérande
            valeurs.empiler(combiner(s[0], a, b));
        } else if (estNombre(s) || estIdentificateur(s)) {
            valeurs.empiler(feuille(s));
        } else {
            throw std::invalid_argument("Jeton invalide '" + s + "'");
        }
    }
    if (valeurs.taille() != 1)
        throw std::invalid_argument("Trop d'operandes (operateur manquant)");
    return valeurs.depiler();
}

// Préfixe -> infixe (entièrement parenthésée)
static std::string prefixeVersInfixe(const std::string& prefixe) {
    std::string s = parcourirPrefixe<std::string>(prefixe,
        [](const std::string& o) { return o; },
        [](char op, const std::string& a, const std::string& b) {
            return "(" + a + " " + std::string(1, op) + " " + b + ")";
        });
    // les parenthèses extérieures sont superflues
    if (s.size() > 1 && s.front() == '(' && s.back() == ')')
        s = s.substr(1, s.size() - 2);
    return s;
}
\end{lstlisting}

\begin{table}[H]
\centering
\small
\begin{tabular}{cll}
  \toprule
  \textbf{Jeton (droite $\to$ gauche)} & \textbf{Action} & \textbf{Pile (sommet à droite)} \\
  \midrule
  \code{c} & empiler & \code{c} \\
  \code{b} & empiler & \code{c | b} \\
  \code{a} & empiler & \code{c | b | a} \\
  \code{+} & $a \gets$ \code{a}, $b \gets$ \code{b} & \code{c | (a + b)} \\
  \code{*} & $a \gets$ \code{(a + b)}, $b \gets$ \code{c} & \code{((a + b) * c)} \\
  \bottomrule
\end{tabular}
\caption{Trace de la conversion de \code{* + a b c} en infixe :
résultat \code{(a + b) * c}}
\end{table}

\subsection{Résultats}

\begin{table}[H]
\centering
\small
\begin{tabular}{>{\ttfamily}l >{\ttfamily}l >{\ttfamily}l}
  \toprule
  \normalfont\textbf{Infixe} & \normalfont\textbf{Préfixe} & \normalfont\textbf{Postfixe} \\
  \midrule
  a + b                        & + a b             & a b + \\
  a + b * c                    & + a * b c         & a b c * + \\
  (a + b) * c                  & * + a b c         & a b + c * \\
  a - b - c                    & - - a b c         & a b - c - \\
  a \textasciicircum{} b \textasciicircum{} c & \textasciicircum{} a \textasciicircum{} b c & a b c \textasciicircum{} \textasciicircum{} \\
  \{ [ (a + b) * c ] - d \} / e & / - * + a b c d e & a b + c * d - e / \\
  (x - y) * (x + y) \textasciicircum{} 2 & * - x y \textasciicircum{} + x y 2 & x y - x y + 2 \textasciicircum{} * \\
  \bottomrule
\end{tabular}
\caption{Conversion infixe $\to$ préfixe et postfixe}
\end{table}

\begin{table}[H]
\centering
\small
\begin{tabular}{>{\ttfamily}l l}
  \toprule
  \normalfont\textbf{Préfixe} & \textbf{Infixe obtenue} \\
  \midrule
  + a b               & \code{a + b} \\
  * + a b c           & \code{(a + b) * c} \\
  - - a b c           & \code{(a - b) - c} \\
  - a - b c           & \code{a - (b - c)} \\
  / - * + a b c d e   & \code{(((a + b) * c) - d) / e} \\
  \midrule
  + a                 & erreur : opérandes manquants pour \code{+} \\
  + a b c             & erreur : trop d'opérandes (opérateur manquant) \\
  + a \$              & erreur : jeton invalide \code{\$} \\
  \bottomrule
\end{tabular}
\caption{Conversion préfixe $\to$ infixe}
\end{table}

\paragraph{Vérification croisée.} Pour valider les conversions, chaque
expression numérique est évaluée de trois manières : directement en
infixe, via sa forme préfixe, et après un aller-retour
infixe $\to$ préfixe $\to$ infixe. Les trois valeurs doivent coïncider.

\begin{lstlisting}[style=sortie, caption={Extrait de la sortie de \code{test\_conversion}}]
  2 ^ 3 ^ 2
     prefixe = ^ 2 ^ 3 2
     retour  = 2 ^ (3 ^ 2)
     valeurs = 512 / 512 / 512   [OK]
  (8 - 2) * (8 + 2) ^ 2
     prefixe = * - 8 2 ^ + 8 2 2
     retour  = (8 - 2) * ((8 + 2) ^ 2)
     valeurs = 600 / 600 / 600   [OK]

6/6 tests reussis
\end{lstlisting}

Les six tests croisés réussissent : la conversion préserve la sémantique de
l'expression, y compris la priorité des opérateurs et l'associativité à
gauche ($-$, $/$) comme à droite (\code{\textasciicircum}). Les expressions préfixes mal
formées sont rejetées avec un message explicite.

% ===========================================================================
\section{Bibliothèque dynamique \texorpdfstring{\code{libpile.so}}{libpile.so} sous Linux}
\label{sec:bibliotheque}

\subsection{Problématique}

Les piles et l'analyseur doivent être fournis sous forme de
\textbf{bibliothèque dynamique} : du code compilé une seule fois, chargé au
lancement par tous les programmes qui l'utilisent. Sous Windows, il s'agit
d'une DLL (\code{.dll}) ; sous Linux, système utilisé pour ce TP, d'un
\emph{shared object} (\code{.so}).

La difficulté vient des \textbf{templates}. \code{PileTableau<T>} n'est
qu'un modèle : le compilateur ne produit du code machine que lorsqu'il est
utilisé avec un type précis (\code{PileTableau<int>}), c'est
l'\emph{instanciation}. Dans la première version du TP, tout le code était
dans les en-têtes (\emph{header-only}) : une bibliothèque compilée à partir
de ces fichiers aurait été \textbf{vide}, puisqu'aucun type n'y est
instancié.

\subsection{Méthodologie}

\subsubsection{Séparation déclaration / implémentation}

Chaque classe générique est répartie sur trois fichiers :
\begin{center}
\small
\begin{tabularx}{\linewidth}{l X}
  \toprule
  \textbf{Fichier} & \textbf{Contenu} \\
  \midrule
  \code{PileTableau.h}   & déclaration de la classe et liste des types déjà
                           compilés dans la bibliothèque (\code{extern template}) \\
  \code{PileTableau.tpp} & code des méthodes du template, écrit hors de la classe \\
  \code{PileTableau.cpp} & instanciation explicite pour les types utilisés,
                           compilé dans \code{libpile.so} \\
  \bottomrule
\end{tabularx}
\end{center}

L'extension \code{.tpp} (\emph{template implementation}) est une
convention : c'est un fichier destiné à être \textbf{inclus}, jamais
compilé seul, ce qui le distingue d'un \code{.cpp}.

\begin{lstlisting}[caption={Instanciation explicite (\code{src/PileTableau.cpp}, compilé dans la bibliothèque)}]
#include "PileTableau.tpp"

template class PileTableau<int>;           // génère ici tout le code de PileTableau<int>
template class PileTableau<double>;
template class PileTableau<char>;
template class PileTableau<std::string>;
\end{lstlisting}

\begin{lstlisting}[caption={Déclaration vue par les programmes clients (fin de \code{include/PileTableau.h})}]
// Instanciations fournies par libpile.so : le compilateur ne les régénère
// pas dans le programme client, il utilise celles de la bibliothèque.
extern template class PileTableau<int>;
extern template class PileTableau<double>;
extern template class PileTableau<char>;
extern template class PileTableau<std::string>;
\end{lstlisting}

Le même principe s'applique à \code{PileListe<T>} et à l'analyseur :
\code{src/Expression.cpp} contient les fonctions utilitaires
(\code{decouper}, \code{appliquer}, \code{priorite}\dots) et instancie
\code{PileTableau<Token>}, \code{PileListe<Token>},
\code{Analyseur<PileTableau>} et \code{Analyseur<PileListe>}.

Une bibliothèque ne peut pas connaître tous les types possibles. Pour le
type utilisateur \code{Point} de l'exercice 1, \code{test\_pile.cpp}
inclut \code{PileTableau.tpp} : \code{PileTableau<Point>} est alors généré
dans le programme lui-même, tandis que \code{PileTableau<int>} provient
toujours de la bibliothèque.

\subsubsection{Construction de la bibliothèque}

\begin{lstlisting}[style=sortie, caption={Étapes de construction (automatisées par le \code{Makefile})}]
# 1. objets en code independant de la position (-fPIC)
g++ -Wall -Wextra -O2 -Iinclude -fPIC -c src/PileTableau.cpp -o obj/PileTableau.o
g++ -Wall -Wextra -O2 -Iinclude -fPIC -c src/PileListe.cpp   -o obj/PileListe.o
g++ -Wall -Wextra -O2 -Iinclude -fPIC -c src/Expression.cpp  -o obj/Expression.o

# 2. regroupement des objets dans la bibliotheque partagee
g++ -shared -Wl,-soname,libpile.so obj/*.o -o libpile.so

# 3. edition de liens d'un programme client avec libpile.so
g++ -Wall -Wextra -O2 -Iinclude src/test_pile.cpp -L. -lpile -Wl,-rpath,'$ORIGIN' -o test_pile
\end{lstlisting}

\begin{itemize}
  \item \code{-fPIC} : une \code{.so} est chargée à une adresse différente
        à chaque exécution, son code ne doit pas dépendre d'une adresse fixe ;
  \item \code{-shared} : produit une bibliothèque partagée au lieu d'un
        exécutable ;
  \item \code{-L. -lpile} : cherche \code{libpile.so} dans le dossier
        courant ;
  \item \code{-Wl,-rpath,'\$ORIGIN'} : enregistre dans l'exécutable qu'il
        doit chercher la bibliothèque dans son propre dossier, sans avoir à
        définir \code{LD\_LIBRARY\_PATH}.
\end{itemize}

\subsection{Résultats}

Les trois programmes de test donnent exactement les mêmes sorties qu'avec la
version \emph{header-only}. Les outils Linux confirment qu'ils utilisent bien
la bibliothèque :

\begin{lstlisting}[style=sortie, caption={Vérification de la liaison dynamique}]
$ ldd test_pile | grep pile
    libpile.so => .../tp1/libpile.so

$ readelf -d test_expression | grep -E 'NEEDED|RUNPATH'
    (NEEDED)   Shared library: [libpile.so]
    (RUNPATH)  Library runpath: [$ORIGIN]

$ nm -DC --undefined-only test_expression | grep Analyseur
    U Analyseur<PileTableau>::verifier(std::string const&)
    U Analyseur<PileListe>::evaluer(std::string const&)
    U Analyseur<PileListe>::verifier(std::string const&)
\end{lstlisting}

Le symbole \code{U} (\emph{undefined}) indique que \code{test\_expression}
ne contient pas le code de l'analyseur : il est résolu au lancement dans
\code{libpile.so}, qui exporte 336 symboles de piles et d'analyseur.

\paragraph{Comparaison avec une bibliothèque statique.} Les mêmes objets
peuvent être archivés en bibliothèque statique (\code{ar rcs libpile.a
obj/*.o}), dont le code est \emph{copié} dans l'exécutable à l'édition de
liens.

\begin{table}[H]
\centering
\small
\begin{tabularx}{\linewidth}{l X X}
  \toprule
  & \textbf{Dynamique} (\code{libpile.so}) & \textbf{Statique} (\code{libpile.a}) \\
  \midrule
  Taille de \code{test\_expression} & 27 Ko & 175 Ko \\
  Bibliothèque absente au lancement & erreur : \code{cannot open shared object file} & fonctionne \\
  Mise à jour de la bibliothèque & remplacer le \code{.so} suffit & refaire l'édition de liens \\
  Mémoire & une copie partagée entre les programmes & une copie par programme \\
  \bottomrule
\end{tabularx}
\caption{Bibliothèque dynamique et statique (mesures sur le TP)}
\end{table}

\subsection{Étapes de compilation}

Le dossier \code{compilation\_steps/} contient un programme de démonstration
(includes, macros, compilation conditionnelle, commentaires, espaces) qui
utilise \code{libpile.so}. Son \code{Makefile} exécute chaque étape
séparément et conserve le fichier produit :

\begin{table}[H]
\centering
\small
\begin{tabular}{llll}
  \toprule
  \textbf{Étape} & \textbf{Commande} & \textbf{Fichier produit} & \textbf{Observation} \\
  \midrule
  prétraitement   & \code{g++ -E} & \code{programme.ii} & 37\,023 lignes, macros remplacées \\
  compilation     & \code{g++ -S} & \code{programme.s}  & code assembleur x86-64 \\
  assemblage      & \code{g++ -c} & \code{programme.o}  & méthodes de la pile non définies (\code{U}) \\
  édition de liens & \code{g++ -lpile} & \code{programme} & lié à \code{libpile.so} (\code{ldd}) \\
  \bottomrule
\end{tabular}
\caption{Fichiers générés dans \code{compilation\_steps/}}
\end{table}

{\sloppy
Après le prétraitement, les commentaires ont disparu, les \code{\#include}
sont remplacés par le contenu des en-têtes (d'où les 37\,023 lignes, dues
surtout à \code{<iostream>}), \code{CARRE(i)} devient \code{((i) * (i))} et
le bloc \code{\#ifdef DEBUG} est supprimé. Le fichier objet
\code{programme.o} référence \code{PileTableau<int>::empiler} sans le
définir : c'est l'édition de liens qui le relie à \code{libpile.so}.\par}

% ===========================================================================
\section{Conclusion}

Ce TP a permis de concevoir une \textbf{pile générique} réutilisable et de
montrer son intérêt sur un problème concret d'analyse d'expressions.

\begin{itemize}
  \item La séparation entre l'interface \code{Pile<T>} et ses deux
        implémentations permet de changer de représentation sans toucher au
        code client : les analyseurs \code{Analyseur<PileTableau>} et
        \code{Analyseur<PileListe>} donnent des résultats identiques.
  \item Le \textbf{tableau dynamique} est environ dix fois plus rapide en
        pratique grâce à la contiguïté mémoire et au doublement de capacité ;
        la \textbf{liste chaînée} garantit un $O(1)$ strict par opération.
  \item La pile rend la vérification des délimiteurs simple et linéaire
        ($O(n)$) ; couplée à un automate à deux états, elle détecte toutes
        les erreurs de syntaxe et les localise précisément.
  \item Un unique algorithme générique à deux piles, paramétré par le type
        des valeurs et une fonction de combinaison, réalise à la fois
        l'évaluation et les conversions infixe $\to$ préfixe/postfixe.
  \item Grâce à l'instanciation explicite des templates et à la séparation
        \code{.h} / \code{.tpp}, le code est distribué sous forme de
        bibliothèque dynamique \code{libpile.so} sous Linux : les programmes
        clients sont plus petits et profitent d'une mise à jour de la
        bibliothèque sans être recompilés.
\end{itemize}

\paragraph{Limites et perspectives.} Les opérateurs unaires (par exemple
$-3$) et les fonctions ($\sin$, $\max$\dots) ne sont pas pris en charge ; ils
pourraient être ajoutés en enrichissant l'automate d'un état et la table des
priorités. L'évaluation pourrait également accepter un environnement
associant une valeur à chaque variable.

% ===========================================================================
\appendix
\section{Compilation et exécution : le Makefile}

\subsection{Rôle d'un Makefile}

Un \code{Makefile} est un fichier de \textbf{recettes de compilation} lu par
la commande \code{make}. Il ne remplace pas l'exécutable : il décrit
\emph{comment le fabriquer}. Sans lui, construire le TP demanderait de taper
à chaque modification une dizaine de commandes \code{g++} longues, dans le
bon ordre et sans oublier d'option (\code{-fPIC}, \code{-rpath}\dots).

Son intérêt principal est la \textbf{compilation incrémentale} :
\code{make} compare la date de modification de chaque fichier avec celle de
ses dépendances et ne relance que les commandes nécessaires. Si seul
\code{src/test\_pile.cpp} est modifié, la bibliothèque n'est pas recompilée ;
si un en-tête change, tout ce qui en dépend est reconstruit.

\subsection{Syntaxe d'une règle}

\begin{lstlisting}[language=make, style=sortie, caption={Forme générale d'une règle}]
cible: dependances
	commande qui fabrique la cible a partir des dependances
\end{lstlisting}

\begin{itemize}
  \item la \textbf{cible} est le fichier à produire (ou un nom d'action comme
        \code{clean}) ;
  \item les \textbf{dépendances} sont les fichiers nécessaires : si l'une
        d'elles est plus récente que la cible, la commande est relancée ;
  \item la \textbf{commande} (la \og recette \fg{}) doit commencer par une
        \textbf{tabulation}, pas par des espaces.
\end{itemize}

Le Makefile du TP utilise aussi des \textbf{variables}
(\code{CXX = g++}, utilisée sous la forme \code{\$(CXX)}), des
\textbf{règles génériques} avec \code{\%} (une seule règle pour tous les
fichiers \code{src/\%.cpp}) et des \textbf{variables automatiques} :

\begin{center}
\small
\begin{tabular}{ll}
  \toprule
  \textbf{Variable} & \textbf{Signification} \\
  \midrule
  \code{\$@} & la cible de la règle \\
  \code{\$<} & la première dépendance \\
  \code{\$\textasciicircum{}} & toutes les dépendances \\
  \code{| obj} & dépendance \og d'ordre \fg{} : le dossier doit exister,
                 sa date est ignorée \\
  \code{.PHONY} & cibles qui ne sont pas des fichiers (\code{all}, \code{run}, \code{clean}\dots) \\
  \bottomrule
\end{tabular}
\end{center}

\subsection{Construction de la bibliothèque et des tests}

\begin{lstlisting}[language=make, style=sortie, caption={\code{Makefile} (1/2) : bibliotheque et programmes de test}]
CXX      = g++
CXXFLAGS = -Wall -Wextra -O2 -Iinclude
HEADERS  = $(wildcard include/*.h include/*.tpp)
PROGS    = test_pile test_expression test_conversion

# Bibliotheque partagee (equivalent Linux d'une DLL)
LIB      = libpile.so
LIB_OBJS = obj/PileTableau.o obj/PileListe.o obj/Expression.o

all: $(LIB) $(PROGS)

# -fPIC : code independant de la position, obligatoire dans une .so
obj/%.o: src/%.cpp $(HEADERS) | obj
	$(CXX) $(CXXFLAGS) -fPIC -c $< -o $@

$(LIB): $(LIB_OBJS)
	$(CXX) -shared -Wl,-soname,$(LIB) $^ -o $@

# programmes lies a libpile.so, cherchee dans leur propre dossier
%: src/%.cpp $(HEADERS) $(LIB)
	$(CXX) $(CXXFLAGS) $< -L. -lpile -Wl,-rpath,'$$ORIGIN' -o $@

obj:
	mkdir -p obj

run: all
	./test_pile && ./test_expression && ./test_conversion
\end{lstlisting}

\begin{center}
\small
\begin{tabularx}{\linewidth}{l X}
  \toprule
  \textbf{Élément} & \textbf{Rôle} \\
  \midrule
  \code{CXX}, \code{CXXFLAGS} & compilateur et options communes : avertissements
      (\code{-Wall -Wextra}), optimisation (\code{-O2}), dossier des en-têtes
      (\code{-Iinclude}) \\
  \code{HEADERS} & liste des \code{.h} et \code{.tpp} (fonction \code{wildcard}) :
      modifier un en-tête déclenche la recompilation de ce qui en dépend \\
  \code{all} & cible par défaut (lancée par \code{make} seul) : la bibliothèque
      puis les trois programmes \\
  \code{obj/\%.o} & compile chaque \code{src/X.cpp} de la bibliothèque en
      \code{obj/X.o}, avec \code{-fPIC} \\
  \code{\$(LIB)} & regroupe les trois objets dans \code{libpile.so}
      (\code{-shared}) \\
  \code{\%} (programmes) & compile \code{src/test\_X.cpp} en exécutable lié à
      \code{libpile.so} ; \code{\$\$ORIGIN} est écrit \code{\$\$} car
      \code{\$} est un caractère spécial pour \code{make} \\
  \code{obj} & crée le dossier des fichiers objets \\
  \code{run} & construit tout puis exécute les trois séries de tests \\
  \bottomrule
\end{tabularx}
\end{center}

\subsection{Cible \texorpdfstring{\code{etapes}}{etapes} : un fichier par étape de compilation}

\begin{lstlisting}[language=make, style=sortie, caption={\code{Makefile} (2/2) : etapes de compilation et nettoyage}]
ETAPES = compilation

etapes: $(ETAPES)/test_pile.ii $(ETAPES)/test_pile.s $(ETAPES)/test_pile.o \
        $(ETAPES)/libpile.a $(ETAPES)/libpile.so \
        $(ETAPES)/test_pile $(ETAPES)/test_pile_statique

# 1. pretraitement : includes recopies, macros remplacees, commentaires supprimes
$(ETAPES)/%.ii: src/%.cpp $(HEADERS) | $(ETAPES)
	$(CXX) $(CXXFLAGS) -E $< -o $@

# 2. compilation : C++ pretraite -> assembleur x86-64
$(ETAPES)/%.s: $(ETAPES)/%.ii
	$(CXX) $(CXXFLAGS) -S $< -o $@

# 3. assemblage : assembleur -> code objet (non executable)
$(ETAPES)/%.o: $(ETAPES)/%.s
	$(CXX) -c $< -o $@

# 4. bibliotheques : archive statique (.a) et bibliotheque partagee (.so)
$(ETAPES)/libpile.a: $(LIB_OBJS) | $(ETAPES)
	ar rcs $@ $^

$(ETAPES)/libpile.so: $(LIB) | $(ETAPES)
	cp $< $@

# 5. edition de liens : dynamique (libpile.so chargee au lancement)
#    et statique (code de libpile.a copie dans l'executable)
$(ETAPES)/test_pile: $(ETAPES)/test_pile.o $(ETAPES)/libpile.so
	$(CXX) $< -L$(ETAPES) -lpile -Wl,-rpath,'$$ORIGIN' -o $@

$(ETAPES)/test_pile_statique: $(ETAPES)/test_pile.o $(ETAPES)/libpile.a
	$(CXX) $^ -o $@

$(ETAPES):
	mkdir -p $(ETAPES)

clean:
	rm -rf obj $(LIB) $(PROGS) $(ETAPES)

.PHONY: all run etapes clean
\end{lstlisting}

La cible \code{etapes} enchaîne les étapes de la compilation de
\code{test\_pile.cpp} en conservant le fichier produit par chacune dans le
dossier \code{compilation/}. Chaque règle prend en entrée le fichier de
l'étape précédente : \code{make} déduit seul l'ordre d'exécution.

\begin{table}[H]
\centering
\small
\begin{tabularx}{\linewidth}{l l r >{\raggedright\arraybackslash}X}
  \toprule
  \textbf{Fichier} & \textbf{Commande} & \textbf{Taille} & \textbf{Contenu} \\
  \midrule
  \code{test\_pile.ii} & \code{g++ -E} & 1\,084 Ko & code prétraité : 42\,160 lignes
      (en-têtes recopiés, macros remplacées) \\
  \code{test\_pile.s} & \code{g++ -S} & 48 Ko & assembleur x86-64 (2\,159 lignes) \\
  \code{test\_pile.o} & \code{g++ -c} & 29 Ko & code objet ; méthodes de la pile
      non définies \\
  \code{libpile.a} & \code{ar rcs} & 328 Ko & bibliothèque statique (archive des
      trois \code{.o}) \\
  \code{libpile.so} & \code{g++ -shared} & 210 Ko & bibliothèque dynamique \\
  \code{test\_pile} & \code{g++ -lpile} & 29 Ko & exécutable lié dynamiquement :
      charge \code{libpile.so} au lancement \\
  \code{test\_pile\_statique} & \code{g++ libpile.a} & 58 Ko & exécutable autonome :
      le code utile de \code{libpile.a} est copié dedans \\
  \bottomrule
\end{tabularx}
\caption{Fichiers générés par \code{make etapes}}
\end{table}

L'exécutable statique est deux fois plus gros que l'exécutable dynamique,
mais bien plus petit que \code{libpile.a} : l'éditeur de liens ne copie que
les fichiers objets de l'archive dont le programme a besoin (les piles, pas
l'analyseur d'expressions).

\subsection{Utilisation}

\begin{lstlisting}[style=sortie]
$ make          # compile libpile.so puis les trois programmes
$ make run      # execute les trois series de tests
$ make etapes   # genere un fichier par etape dans compilation/
$ make clean    # supprime tous les fichiers generes
\end{lstlisting}

\end{document}
